\section{Introduction}

With the rapid development of the Internet, online recruitment platforms (e.g., LinkedIn\footnote{https://www.linkedin.com}{}, Indeed\footnote{https://www.indeed.com}{}, and ZipRecruiter\footnote{https://www.ziprecruiter.com}{}) are becoming essential for recruiting and job seeking. A considerable number of talent profiles and job descriptions are posted on these platforms. Taking LinkedIn as an example, more than 900 million members have registered, and 90 jobs were posted every second by the first quarter of 2023\footnote{https://news.linkedin.com/about-us\#Statistics}{}. Considering such a large number of options, Person-Job Fit (PJF) \cite{malinowski2006matching} has become a critical research topic for improving the efficiency of recruitment and job seeking. Person-Job Fit aims to automatically link the right talents to the right job positions according to talent competencies and job requirements. 

Previous works on Person-Job Fit mainly focus on leveraging two types of information, namely, (1) historical job application information and (2) textual information in profiles and job descriptions. Collaborative filtering-based methods \cite{shalaby2017help,bian2020learning} are applied to capture co-apply relations between job seekers and co-applied relations between job positions in the historical application information. Manually-engineered textual features and deep language models have been widely adopted to leverage the textual information \cite{bian2019domain,qin2018enhancing, yan2019interview, lu2013recommender}.

However, workplace social connections among members, commonly called professional networks, have been overlooked as a pivotal source of information. A survey reported on the Official LinkedIn Blog\footnote{\url{https://blog.linkedin.com/2017/june/22/the-best-way-to-land-your-next-job-opportunity}} showed that out of more than 15,000 LinkedIn members, 80\% believed that professional networking could help find new job opportunities, an even more striking 70\% gained job opportunities directly through their connections. As an important source for achieving employment relationships, professional networks have great potential for helping with Person-Job Fit. Incorporating professional networks into the models can offer two main benefits: (1) Professional networks directly improve Person-Job Fit by bridging the gap between job seekers and potential job opportunities. A Person-Job Fit model enriched with professional network information can suggest job opportunities to users based on their professional connections. (2) Professional networks can also help complete job seekers' profiles. Often job seekers' online profiles lack comprehensive details, yet one's professional experience and skills can be, to some extent, discernible through their professional connections.

While professional networks offer advantages, they often contain a lot of noise. This noise includes connections that aren't relevant and information from relevant connections that don't help improve Person-Job Fit. To illustrate, irrelevant connections might involve job seekers' former classmates or recruiters from unrelated industries. Even when the connections are relevant to the job seeker, they could still have information unrelated to the job seeker's aspirations for the role. To tackle these noises, a promising strategy involves using heterogeneous knowledge, such as job seekers' skills, work experience, and educational background, to gauge professional connections' relevance and extract pertinent information to elevate Person-Job Fit.

In this paper, we propose a graph neural network-based framework that utilizes heterogeneous knowledge to integrate professional networks into Person-Job Fit. We address the challenge of social noise in professional networks by designing a job-specific attention mechanism. Initially, we define a Workplace Heterogeneous Information Network (WHIN) that captures heterogeneous knowledge, including professional connections. We employ a heterogeneous graph pre-training technique to learn the representations of various entities in the WHIN. Subsequently, we introduce the Contextual Social Attention Graph Neural Network (CSAGNN), designed to supplement users' lacking information with contextual information from their professional connections. To tackle the social noise in the workplace social network, we infuse a job-specific attention mechanism into CSAGNN, capitalizing on the pre-trained entity representations from WHIN.

The main contributions of this study are as follows:
\begin{itemize}
    \item To our knowledge, we are the first to define a heterogeneous information network that incorporates heterogeneous knowledge in the Person-Job Fit scenario.
    
    \item We systematically utilize professional networks in a two-stage approach, WHIN pre-training and CSAGNN, to address the Person-Job Fit task.
    
    \item We present a novel Contextual Social Attention Graph Neural Network (CSAGNN) specifically designed to handle noisy professional networks, effectively mitigating the impact of irrelevant information while focusing on related professional connections and contexts for Person-Job Fit.
    
    \item We evaluate our approach on three real-world datasets across diverse industries. Experimental results show that our model outperforms baseline models.
\end{itemize}
